\scanpage{12}\mspage{9}

Alors on voit que
\[
 L^{n-i+\alpha}q_{i\alpha}: M^i\unsure{-|\to}M^{2n-i+2\alpha}
\]
est nul sur les $\deleted{\unsure{P}}L^{\alpha'}P^{i-2\alpha'}$ pour $\alpha'\ne\alpha$, et que sur $L^\alpha P^{i-2\alpha}$, \deleted{\unsure{c'est}} son composé avec $P^{i-2\alpha}\xrightarrow{L^\alpha} L^\alpha P^{i-2\alpha}$ n'est autre que l'isom
\[
 L^{n-i-2\alpha}:P^{i-2\alpha}\longrightarrow L^{n-i+2\alpha}P^{i-2\alpha}\subset M^{2n-i+2\alpha}.
\]
Il s'ensuit aussitôt qu'on a
\[
 \pi_{i,\alpha}=L^\alpha w_{i-2\alpha}L^{n-i+\alpha}q_{i\alpha}
 =L^\alpha w_{i-2\alpha}L^{n-i+\alpha}(1-\sum_{\alpha'>\alpha}\pi_{i,\alpha'}\unsure{)}
 \tag*{(6.9)}
\]
ce qui donne bien $\pi_{i\alpha}\in\unsure{E|\mathcal E}$. \deleted{Cela prouve (iv).} Si $i<n$, il faut \unsure{égalt} $\pi_{2n-i,\alpha}\in E$. On

\deleted{Reste à prouver \inserted{(iii)} $\Lambda\in E$. Il suffit de \deleted{\unsure{prouv.}} \unsure{poser}
\[
 q_{2n-i,\alpha}=1-\sum_{\alpha'>\alpha}\pi_{2n-i,\alpha'}
\]
et on voit alors que
\[
 \deleted{\unsure{q_{2n-i,\alpha}|\pi_{2n-i,\alpha}}} L^\alpha q_{2n-i,\alpha}:M^{2n-i}\unsure{-|\to}M^{2n-i+2\alpha}
\]}
le déduit \unsure{p\textsuperscript{r} ex.|par ex.} de la formule
\[
 \pi_{2n-i,\alpha}=L^{n-i}\pi_{i,\alpha}w_i\deleted{\unsure{\pi_{2n-i}}}
 \tag*{(6.10)}
\]
Cela prouve \inserted{(6.8') donc} (iv). Reste à prouver (iii) $\Lambda\in E$.

Comme on a
\[
 \Lambda=\sum_{\substack{j,\alpha\\\deleted{\unsure{0\le j\le 2n}}}}\deleted{\unsure{\{}}\Lambda\pi_{j,\alpha}
 \deleted{\unsure{+\sum_{\substack{\alpha>0\\0\le i\le n}}\Lambda\pi_{i,\alpha}\}}}
 \tag*{(6.11)}
\]
il suffit de prouver que $\Lambda\pi_{i\alpha}\deleted{\unsure{\Lambda\pi_{2n-i,\alpha}}}\in E$.
% archive sheet number: 11
